% TOP_PROBLEM: {"id": 2564, "title": "Kourovka Notebook Problem 21.55", "category_id": 17, "category": {"id": 17, "name": "group_theory", "display_name": "Group Theory", "description": "Problems about groups, group actions, representations, and related algebraic structures.", "slug": "group-theory", "order_index": 17, "created_at": "2026-07-31T15:26:25.671Z"}, "status": "open", "problem_number": "KOU-21.55", "set_id": 10}
% FIRST_POSTED: 2026-10-01
% ENTRY_KIND: statement_only
% UnsolvedMath ID 2564; source catalogue code KOU-21.55
% !TeX program = lualatex
% Definition and sources only; this file is not an LLM solution attempt.
\documentclass[11pt,a4paper]{article}
\usepackage[margin=25mm]{geometry}
\usepackage{fontspec}
\setmainfont{lmroman10-regular.otf}[BoldFont=lmroman10-bold.otf,
 ItalicFont=lmroman10-italic.otf,BoldItalicFont=lmroman10-bolditalic.otf]
\usepackage{amsmath,amssymb,mathtools}
\usepackage{unicode-math}
\setmathfont{latinmodern-math.otf}
\AtBeginDocument{\let\setminus\smallsetminus}
\usepackage{xurl,hyperref}
\hypersetup{colorlinks=true,urlcolor=blue}
\setcounter{secnumdepth}{0}
\setlength{\parindent}{0pt}
\setlength{\parskip}{5pt}
\setlength{\emergencystretch}{3em}
\begin{document}
\section*{Kourovka Notebook Problem 21.55}\label{2564}
{\small UnsolvedMath ID 2564 · KOU-21.55 · Group Theory · Kourovka Notebook - New Problems, Issue 21}
\subsection{Definitions and mathematical statement}
Fix a prime power $q=p^a$, with $p$ prime and $a\ge1$, and let
$\mathbb F_q$ be the field with $q$ elements.
A finite group $H$ is $p$-soluble if it has a finite normal series
whose factors are either $p$-groups or groups of order prime to $p$.
The $p$-length $\ell_p(H)$ is the minimum number of nontrivial
$p$-group factors among such series. The factors of order prime to
$p$ do not contribute to this count. In particular a group of order
prime to $p$ has $p$-length zero.

For each positive integer $n$, define
\[
 m_n(q)=\max\{\ell_p(H):H\le\operatorname{GL}_n(\mathbb F_q),
                                      \ H\text{ is }p\text{-soluble}\}.
\]
The maximum exists because the ambient general linear group is finite.
Subgroups need not act irreducibly, and are not required to be soluble
at primes other than the specified $p$.

The question is whether, for every fixed $q$,
\[
 \lim_{n\to\infty}\frac{m_n(q)}{\log_2 n}=1.
\]
Equivalently, for every $\varepsilon>0$ there should be $N(q,\varepsilon)$
such that for all $n\ge N(q,\varepsilon)$ the ratio differs from one
by less than $\varepsilon$.

The field is fixed before the dimension tends to infinity; uniformity
as $q$ varies is not requested. The logarithm is specifically base two,
so its choice fixes the leading constant. Both an asymptotic upper
bound for all eligible subgroups and a matching lower construction
are needed for the asserted limit.
\subsection{Short English statement}
For a fixed finite field of characteristic $p$, is the maximal $p$-length of a $p$-soluble linear subgroup asymptotic to the base-two logarithm of the dimension?
\subsection{Sources}
\begin{itemize}
\item[S1] ULAM.AI and the UnsolvedMath Contributors, supplied problems.json, record 2564 (KOU-21.55), Kourovka Notebook - New Problems, Issue 21. Catalogue identity and source transcription; CC BY 4.0.
\url{https://huggingface.co/datasets/ulamai/UnsolvedMath}.
\item[S2] The Kourovka Notebook, 21st edition, new problems, Problem 21.55. Expanded formulation of the supplied catalogue statement.
\url{https://arxiv.org/pdf/1401.0300v47}.
\end{itemize}
\end{document}
